\section{河流、疾病与迁居：不在帝纪中心的金代北方}

\subsection{渡口决定逃与留}

北方战争中的迁徙不能只由地图上的箭头表示。一个家庭能否离开，取决于渡口是否开放、桥梁是否仍可过车、是否有亲属接应、携带的粮食能支持几日，以及下一处城镇是否允许陌生人停留。黄河、淮河及其支流既是运输与防御的条件，也是逃亡和重返家园的条件。河水上涨、渡船被军队征用或城门实行检查，都可能使原本可行的路径突然消失。

正史所称的流民，常是行政需要下的集合词。它可指失地者、避兵者、被征调后离散者或暂时迁往近城的人；这些人的目的地、财产与亲属关系未必相同。用一个巨大数字描写流离失所很能制造悲剧，却容易掩盖数字来自何种报告、在何地统计、是否包含重复移动。没有足够材料时，本章宁可确认战争与征发增加了迁居压力，而不替无名者安排同一条路线。

\subsection{环境不是命运，却规定了代价}

旱、涝、蝗灾和疫病在金代文字中反复出现。它们既可能记录真实的地方困境，也可能被写入灾异政治：官员以灾情请求蠲免，朝廷以救灾展示德政，史家以天灾评判治乱。因此一条灾异记录不足以推出全境收成或人口变化。考古沉积、河道分期和城址废弃能提供另一类线索，但其年代和空间范围也要逐件核对。

环境的历史意义在于它如何同制度和战争相遇。雨水使一段道路泥泞，军粮到不了前线，地方的车马和劳力便可能被再次征集；河道变化使田地受损，县衙的税额和家户的储粮就发生冲突；疫病进入拥挤城池，围城与迁徙的风险随之上升。这些并不是一条单向因果链。国家的赈济、地方的互助、市场价格和家庭储备都会改变结果，材料能看见哪一环，叙事便只能写到哪一环。

\subsection{重新安置不是回到原位}

金廷在战争和扩张中需要安置军人、迁入者和官署人员。安置涉及土地、住房、供给、编户和治安；对中枢而言是恢复秩序，对当地人则可能意味着新的邻居、新的役法和对水土的重新竞争。移民被送到某地，并不说明他们已经稳定耕作或与旧居民达成关系。反过来，地方接受新来者也未必只是被动服从：婚姻、租佃、市场和宗教网络都可能让陌生人逐渐嵌入。

墓志中偶见的籍贯与迁居，最适合说明一个家族怎样叙述自己的流动。它通常会将迁徙写成仕宦、祖德或避乱的正确选择，很少记录途中失去的亲人、债务和失败者。题记、地方遗址与钱币则让我们在别处看见活动痕迹。它们不能拼成一部完整的迁徙档案，却足以反驳“政权转换后人口自然归位”的想象。

\subsection{边区的家庭有不同的季节}

沿边军镇的家庭要面对驻军、贸易和警报；内地农户更直接面对田赋、市场和道路；东北的军政迁徙又会把部众关系带入新的地方。相同的法令在这些地区产生不同后果，并非因为某处更“先进”或更“落后”，而是水源、交通、军需和既有社会关系不同。边区家户可能从供应军队中取得收入，也可能因马匹、粮草和劳役而失去耕作时间。两种情形可以同时存在。

金朝的边防因此不只是城墙与军队。它依赖家庭能否生产余粮、市场能否提供物资、道路能否在季节变化中通行，也依赖地方官是否能把急迫命令转成可执行安排。当这些条件被蒙古战争逐步破坏，边防的崩解先表现为消息、补给和安置的失灵，最后才在帝纪中显为失城与亡国。
